\question[5]
Ein Online-Shop verwaltet Informationen über Kunden, Bestellungen und Produkte.

\begin{itemize}
\item Jeder Kunde besitzt eine KundenID und einen Namen.
\item Jede Bestellung besitzt eine BestellID und ein Bestelldatum.
\item Jedes Produkt besitzt eine ProduktID, eine Bezeichnung und einen Preis.
\item Ein Kunde kann mehrere Bestellungen aufgeben, jede Bestellung gehört 
jedoch genau zu einem Kunden.
\item Eine Bestellung kann mehrere Produkte enthalten,
 und ein Produkt kann in mehreren Bestellungen enthalten sein.
\end{itemize}

Modelliere den beschriebenen Ausschnitt der Realität als Entity-Relationship-Diagramm.
Kennzeichne die Schlüssel und gib die Kardinalitäten an.


\begin{solutionbox}{7cm}

\includegraphics[width=9cm]{\pfad/cur04c.png}
% Entitaeten:
% \begin{itemize}
% \item Kunde
% \item Bestellung
% \item Produkt
% \end{itemize}

% Attribute:
% \begin{itemize}
% \item Kunde: KundenID (Schluessel), Name
% \item Bestellung: BestellID (Schluessel), Bestelldatum
% \item Produkt: ProduktID (Schluessel), Bezeichnung, Preis
% \end{itemize}

% Beziehungen und Kardinalitaeten:
% \begin{itemize}
% \item Kunde -- gibt auf -- Bestellung: 1:n
% \item Bestellung -- enthaelt -- Produkt: n:m
% \end{itemize}

% Hinweis:
% Die n:m-Beziehung zwischen Bestellung und Produkt kann im relationalen Modell
% spaeter durch eine Zwischentabelle (z. B. Bestellposition) aufgeloest werden.
\end{solutionbox}
